% !TeX root = ../install-latex-guide-en.tex

% \chapter*{前言}
\chapter*{Preface}

% 以``啸行''名义加入 QQ 群 91940767、478023327、640633524 和 200392395 后,
% 经常有群友咨询如何安装 \LaTeX.
% 实际上,
% 用户安装的是 \LaTeX\ 的\textbf{发行版}和相关的\textbf{编辑器}.
% 本手册将介绍在\textbf{不存在其他 \LaTeX\ 发行版 (如 C\TeX\ 套装) 的前提下},
% 在 Windows 11、Ubuntu 24.04 和 macOS 系统中安装
% \TeX~Live (macOS 中介绍 Mac\TeX)、升级宏包、编译简易文档等相关操作,
% 并多以介绍命令行操作为主.
% 有关 MiK\TeX\ 的安装,
% 可以参考 \href{https://camusecao.top/2021-06-16/MiKTeX/}{MiK\TeX\ 的基本使用}.
This manual will introduce the installation of \TeX~Live
(Mac\TeX\ is introduced in macOS), upgrading packages,
compiling simple documents, and other related operations in Windows 11,
Ubuntu 24.04, and macOS, assuming that
\textbf{other \LaTeX\ distributions (such as C\TeX\ Suite) do not exist},
and mainly introduces command line operations.
For the installation of MiK\TeX, please refer to
\href{https://camusecao.top/2021-06-16/MiKTeX/}{\emph{Basic use of \textup{MiK\TeX}}}.

% 本手册还将简要介绍几款常见编辑器的使用方法,
% 其他编辑器如 \href{https://code.visualstudio.com/}{VS Code},
% \href{https://www.vim.org/}{Vim},
% 用户可自行了解它们的使用方法.
% 鉴于 WSL 较为特殊,
% 本手册提供了一份 VS Code 配合
% \href{https://marketplace.visualstudio.com/items?itemName=James-Yu.latex-workshop}{\LaTeX\ Workshop}
% 的自用配置单,
% 仅供参考.
This manual will also briefly introduce the usage of several common editors.
For other editors such as \href{https://code.visualstudio.com/}{VS Code},
\href{https://www.vim.org/}{Vim}, users can learn how to use them by themselves.
Given that WSL is quite special,
this manual provides a self-use configuration list for VS Code with
\href{https://marketplace.visualstudio.com/items?itemName=James-Yu.latex-workshop}{\LaTeX\ Workshop}
for reference only.

% 本手册还额外介绍在线 \LaTeX\ 平台的相关内容,
% 如国际知名的 \href{http://www.overleaf.com/}{Overleaf} 等.
This manual also provides introductions about online \LaTeX\ platforms,
Such as the internationally renowned
\href{http://www.overleaf.com/}{Overleaf}, etc.

% 本手册所涉及到的代码需结合上下文说明, 不能简单地复制粘贴.
% 粉色文字都是可点的超链接, 可直接跳转.
% \menu{菜单} 表示软件菜单. \keys{k} 表示键盘按键.
% 建议用户阅读 \href{https://tug.org/texlive/doc/texlive-zh-cn/texlive-zh-cn.pdf}{\textsf{texlive-zh-cn}}
% 和 \href{http://mirrors.ctan.org/info/lshort/chinese/lshort-zh-cn.pdf}{\textsf{lshort-zh-cn}}
% 以更全面地了解基础内容.
The codes involved in this manual need to be explained in context
and cannot be simply copied and pasted.
The pink texts are clickable hyperlinks that can be directly jumped.
\menu{menu} represents the software menu, \keys{k} represents the keyboard keys.
It is recommended that users read
\href{https://tug.org/texlive/doc/texlive-en/texlive-en.pdf}{\textsf{texlive-en}}
and
\href{http://mirrors.ctan.org/info/lshort/english/lshort.pdf}{\textsf{lshort}}
to have a more comprehensive understanding of the basic content.

% 必须声明的是,
% 即便你已经通读本手册和上面提到的两本书,
% 有很大概率还会碰到很多问题.
% 这时我希望用户可以按照正确方式提问,
% 例如提供\textbf{最小工作示例},
% 并在\href{https://ask.latexstudio.net/}{论坛}%
% 上面留下你的问题和相应解答以方便后来者.
It must be stated that even if you have read this manual
and the two books mentioned above,
there is a high probability that you will still encounter many questions.
At this time, I hope that users can ask questions in the right way,
such as providing a \textbf{Minimal Working Example},
and leaving your questions and corresponding answers on the
\href{https://tex.stackexchange.com/}{forum}
for the convenience of later generations.

% 本手册大部分内容是个人过去一段时间的使用总结, 其中难免有不甚合理或晦涩难懂的部分.
% 若用户在阅读本手册的过程中有任何意见和建议,
% 请发\href{mailto:ranwang.osbert@outlook.com}{邮件}或在
% \href{https://github.com/OsbertWang/install-latex-guide-zh-cn}{GitHub} 中提
% \href{https://github.com/OsbertWang/install-latex-guide-zh-cn/issues}{issue}.
% 本手册另在\href{https://gitee.com/OsbertWang/install-latex-guide-zh-cn}{码云}%
% 有备份,
% 并于 2020 年 7 月提交至 CTAN.
Most of this manual is a summary of the use of the Chinese author
\href{https://github.com/OsbertWang}{\underline{Ran Wang}}
over the past period of time.
It is inevitable that there are some unreasonable or obscure parts.
If users have any comments and suggestions while reading this manual,
please send \href{mailto:myhsia@outlook.com}{email} or raise
\href{https://github.com/myhsia/install-latex-guide-en}{GitHub}
\href{https://github.com/OsbertWang/install-latex-guide-en/issues}{issue}.
The Chinese version of this manual was submitted to CTAN in July 2020, and the English translation version was submitted to CTAN in July 2025.
\emph{Unlike the Chinese version,
this manual removes some Chinese localization operations.}

% 本手册发布后,
% \href{https://github.com/EthanDeng}{Dongsheng Deng},
% \href{https://github.com/muzimuzhi}{muzimuzhi},
% \href{https://github.com/stone-zeng}{Xiangdong Zeng},
% \href{https://github.com/taumasyang}{tauyoung},
% \href{https://github.com/myhsia}{myhsia},
% 对本手册提出了很好的建议, 并提供了帮助,
% 其中, 有关 macOS 的内容最初由 Xiangdong Zeng 草拟完成,
% 而后 Dongsheng Deng, tauyoung 进行了补充,
% 最近一次更新则由 myhsia 提供.
% 在此一并感谢.
\href{https://github.com/EthanDeng}{Dongsheng Deng},
\href{https://github.com/muzimuzhi}{muzimuzhi},
\href{https://github.com/stone-zeng}{Xiangdong Zeng},
\href{https://github.com/taumasyang}{tauyoung},
\href{https://github.com/myhsia}{myhsia}
have provided good suggestions and help on the
\href{https://github.com/OsbertWang/install-latex-zh-cn}{Chinese version}
of this manual.
The content about macOS was originally drafted by Xiangdong Zeng,
supplemented by Dongsheng Deng, tauyoung,
and the latest update was provided by myhsia.
Thanks to them.

% 本手册自 2024 年 4 月开通捐赠渠道.
% 如果你认为本手册曾经帮助过你,
% 并且你还愿意继续支持本手册,
% 本着自愿原则,
% 你可以扫描以下二维码为本手册捐赠固定金额.
% 需要说明的是,
% 捐赠行为不会影响本手册的更新.
% 衷心感谢你的支持.
\medskip \noindent
\begin{minipage}{.64\linewidth}
  \hspace{15pt}
  To support the development of this manual, welcome scan the QR codes
  on the right to donate a fixed amount for this manual.
  A certain percentage of the donation will be transferred to the author of the
  \href{https://github.com/OsbertWang/install-latex-zh-cn}{Chinese version}
  of this manual.
  Thank you for your support.
\end{minipage}
\hspace*\fill
\begin{minipage}{.32\linewidth}
  \definecolor{paypal}{HTML}{253B80}
  \definecolor{stripe}{HTML}{635bff}
  \begin{center}
    \fancyqr[ classic, height = 5.55em, version = 8, image = {\tikz{
              \node [ inner sep = 0pt, outer sep = 0pt,
                      rounded corners = .225em, draw, paypal,
                      line width = .15em, font = \scriptsize,
                      minimum width = .9em, minimum height = .9em
                    ] (paypal) {\textcolor{paypal}\faCcPaypal};
              \draw [ paypal, very thick ]
                    ([yshift = -.1125em]paypal.north west) --++ (.9em,0)
                    ([yshift = +.1125em]paypal.south west) --++ (.9em,0);}}
            ]{https://www.paypal.com/qrcodes/managed/ 86ffb3de-341c-458f-b410-b4026f90e519?utm_source=consapp_download&  amount=50&currency_code=HKD}%
    \qquad
    \fancyqr[ classic, height = 5.55em, version = 8, image = {\tikz{
              \node [ inner sep = 0pt, rounded corners = .225em,
                      draw, stripe, line width = .15em, font = \scriptsize,
                      minimum width = .9em, minimum height = .9em
                    ] (stripe) {\textcolor{stripe}\faCcStripe};
              \draw [ stripe, very thick ]
                    ([yshift = -.1125em]paypal.north west) --++ (.9em,0)
                    ([yshift = +.1125em]paypal.south west) --++ (.9em,0);}}
            ]{https://donate.stripe.com/6oU14n48obDB8rG2DBbV600}
  \end{center}
\end{minipage}
